\scanpage{141}
$K$ extension finie de $\mathbb Q$
\[
T_K=\prod_{K/\mathbb Q}G_m,\qquad
N_{K/\mathbb Q}:T_K\longrightarrow T_{\mathbb Q}=G_m,\qquad
\deleted{\unsure{V_K}}\quad V_K(\subset T_K)=\operatorname{Ker}N_{K/\mathbb Q};
\]
\[
1\longrightarrow V_K\longrightarrow T_K\longrightarrow G_m\longrightarrow1
\]
\margin{\textcircled{\S\,1}}
\uline{NB} Si $E_K$ unités, alors $E_K\cap V_K(\mathbb Q)$ est d'indice $1$ ou $2$
dans $E_K$, et contient $E_{1K}$ [unités complètes y compris à l'$\infty$]

\uline{Proposition} Soit $V'_{\deleted{\unsure{K}}}$ \inserted{$=V_K/R$} un groupe quotient de $V_K$.
Les conditions suivantes sont équivalentes

(i) L'image de $E_1$ dans $V'_{\deleted{\unsure{K}}}(\mathbb Q)$ est finie, i.e.
$R\cap E$ est d'indice fini dans $E$.

(ii) Il existe un ``module'' $m$ de $K$ tel que
$T_{\deleted{\unsure{K}}}\longrightarrow T_m=T/\overline E_m$
est nul sur $R$, i.e. $T'=T/R_{\deleted{\unsure{E}}}$ est majoré par le quotient $T_m$ de $T$.

(iii) $(V')_{\mathbb R}$ est un tore compact.

(iv) Si $L$ \deleted{\unsure{L}} est le plus-grand sous-corps Weilien de $K$, et
$L_0=L\cap\mathbb R\subset L$ (d'indice $1$ ou $2$ dans $L$)
\deleted{le plus-grand sous-corps tot réel de $K$,}
\deleted{Alors l'hom $i_{K,L}:T_{\deleted{\unsure{K}}L_0}\longrightarrow T_K$
(qui applique $V_K$ dans $V_L$) applique,}
$R$ \deleted{est donc} \inserted{contient} le noyau de l'\deleted{épi}morphisme
\[
\deleted{\unsure{T}}\quad
V_K\xrightarrow{N_{K/L}}V_L\xrightarrow{\mathrm{can}}V_L/\operatorname{Im}V_{L_0}.
\]
où $V_{L_0}\longrightarrow V_L$ est induit par l'inclusion
$T_{L_0}\hookrightarrow T_L$; i.e.
$V'_K$ ``est'' un quotient de $V_L/\operatorname{Im}V_{L_0}$.

\uline{Dém} (i) signifie qu'il existe $F\deleted{E}\subset E_1$
\inserted{ss-gr.} d'indice fini dans $E_1$ telle que $V'$ soit majoré par
$\deleted{T}\inserted{V}/\overline F$ ($\overline F$, adhérence zariskienne de $F$)
i.e. $T'$ majoré par $V/\overline F$; comme les $E_m$ sont d'indice fini
dans $E_1$, on voit que $(\mathrm{ii})\longrightarrow(\mathrm{i})$, et
$(\mathrm{i})\Rightarrow(\mathrm{ii})$ car tout sous-groupe d'indice fini de $E_1$
\deleted{contient} (ou de $E$, c'est kif) contient un $E_m$ (Chevalley).

$(\mathrm{i})\Longrightarrow(\mathrm{iii})$ Car soit \deleted{$E$}
$F=R\cap E$, de sorte que $V'$ est majoré par $V/\overline F$,
\deleted{si} donc $V'_{\mathbb R}$ majoré par $V_{\mathbb R}/\overline F$
(où maintenant $\overline F$ désigne l'adhérence zariskienne dans $V_{\mathbb R}$),
\inserted{et $V'(\mathbb R)=V_{\mathbb R}(\mathbb R)/\overline F(\mathbb R)$.}
Or on sait que
\[
V(\mathbb R)\subset
(\mathbb R^{*r_1}\times\mathbb C^{*r_2})
\inserted{\simeq\mathbb R_+^{*(r_1+r_2)}\times(\mathbb Z/2\mathbb Z)^{r_1}\times U^{r_2}}
\]
[le \deleted{\unsure{diviseur}} sous groupe des syst. $(\lambda_i)(\mu_j)$
tels que $\prod\lambda_i\prod|\mu_j|^2=1$
\inserted{i.e. $(\rho_i)(\varepsilon_j),(u_h)$ tel que $\prod\rho_i=1$}],
et $E_1$ (donc $F$) en est un
\margin{140}
